\cotefile{62}
\scanpage{51}
\mspage{- 27 -}

n\textsuperscript{o} 276

théorème sur les anneaux de Zariski, plutôt que des disperser les diverses étapes de la démonstration en Propositions auxilliaires ainsi qu'il est fait dans l'Etat Serre ; on peut espérer que le lecteur moyen sera bien capable de voir tout seul, par exemple, que la relation $\hat{\mathrm E}=\overset{\handwritten{\scriptscriptstyle\wedge}}{\mathrm A}\handwritten{\otimes}\mathrm E$ est vraie pour tout anneau $\mathrm m$-adique noethérien, vu que la démonstration ne fait pas usage de l'hypothèse $\mathrm m\handwritten{\subset}\mathrm r(\mathrm A)$ \ldots)

% archive sheet number: 50
\scanpage{52}
\uline{Autre façon de procéder}

Prenons la section privilégiée comme section $\infty$, donc $X-s_{\infty}=\Lambda$ apparait comme un $S$-fibré principal \deleted{projectif} \unsure{torseur|homogène} de groupe $\mathcal L$, où $\mathcal L$ est un $S$-schéma en groupes, associé à un \deleted{fibré} faisceau inversible $\underline L$ \margin{$\Vert$} $[\mathord{-}\mathord{-}\mathord{-}\mathord{-}\mathord{-}\unsure{]}$ : expliquer chaque fois pour un schéma de Brauer \unsure{trivisé|trivié|Severi} et une section dans la droite \unsure{--}$]$. Alors la \deleted{\unsure{donnée}} trisection de $\Lambda$ permet de construire, par somme $[\unsure{--}]$ une section de $\Lambda^{3}$, d'où $[3$ \deleted{\unsure{est|étant}} inversible$]$ une section de $\Lambda$, i.e. la \uline{section barycentrique} $[$à expliquer en général, chaque fois que $n$ est inversible sur un schéma en groupes \inserted{\unsure{abéliens}} $\mathcal L$ et qu'on a une $n$-section d'un fibré principal $\Lambda$ de groupe $\mathcal L]$ Cela définit un isomorphisme $\Lambda\simeq\mathcal L$ \uline{global}, identifiant $X'$ à une

% two vertical emphasis strokes beside the bracket, left margin
% archive sheet number: 51
\scanpage{53}
\mspage{- 26 -}

n\textsuperscript{o} 276

on voit que le couple $(\mathrm A,\hat{\mathrm A})$ est plat.

Soient $\mathrm E$ un $\mathrm A$-module de type fini et $\mathrm F$ un sous-module de $\mathrm E$ ; identifiant $\hat{\mathrm E}$ à $\hat{\mathrm A}\handwritten{\otimes}\mathrm E$ on voit que $\hat{\mathrm F}$ est l'image de $\hat{\mathrm A}\handwritten{\otimes}\mathrm F$ dans $\hat{\mathrm E}$ ; autrement dit on a la formule
\[
\hat{\mathrm F}=\hat{\mathrm A}.\mathrm F\quad .
\]

En particulier il vient :
\[
\mathrm m^{\mathrm p}=\hat{\mathrm A}.\mathrm m^{\mathrm p}=(\hat{\mathrm A}.\mathrm m)^{\mathrm p}=(\hat{\mathrm m})^{\mathrm p}\ , 
\]
ce qui montre que la topologie de $\hat{\mathrm A}$ coincide avec la topologie $\hat{\mathrm m}$-adique Comme $\hat{\mathrm A}$ est complet et séparé pour la topologie $\hat{\mathrm m}$-adique, l'idéal $\hat{\mathrm m}$ est contenu dans le radical de $\hat{\mathrm A}$ ; comme $\mathrm G(\hat{\mathrm A})=\mathrm G(\mathrm A)$ est noethérien (parce qu'à engendrement fini sur l'anneau noethérien $\mathrm A/\mathrm m$) on voit d'autre part que $\hat{\mathrm A}$ est noethérien (\handwritten{\S}1, corollaire 2 de la proposition 7). Par suite\unsure{,} le couple $(\hat{\mathrm A},\hat{\mathrm m})$ est un anneau de Zariski, ce qui achève la démonstration.

\uline{Corollaire} - \uline{Soient} $(\mathrm A,\mathrm m)$ \uline{un anneau de Zariski et} $\mathrm E$ \uline{un A-module de type fini}. \uline{Pour tout sous-module} $\mathrm F$ \uline{de} $\mathrm E$ \uline{on a}
\[
\hat{\mathrm F}=\hat{\mathrm A}.\mathrm F\qquad\text{et }\mathrm F=\hat{\mathrm A}.\mathrm F\handwritten{\cap}\mathrm E.
\]

\uline{Si} $\mathrm F$ \uline{et} $\mathrm G$ \uline{sont deux sous-modules de} $\mathrm E$, \uline{on a}
\[
\handwritten{\widehat}{\mathrm F+\mathrm G}=\hat{\mathrm F}+\hat{\mathrm G}\ ;\quad
\handwritten{\widehat}{\mathrm F\handwritten{\cap}\mathrm G}=\hat{\mathrm F}\handwritten{\cap}\hat{\mathrm G}\ ;\quad
(\handwritten{\widehat}{\mathrm F:\mathrm G})=(\hat{\mathrm F}:\hat{\mathrm G}).
\]

La relation $\hat{\mathrm F}=\hat{\mathrm A}.\mathrm F$ a déjà été démontrée ; la relation $\mathrm F=\hat{\mathrm A}.\mathrm F\handwritten{\cap}\mathrm E$ exprime que $\mathrm F$ est fermé dans $\mathrm E$ (théorème 2) : la relation $\handwritten{\widehat}{\mathrm F+\mathrm G}=\hat{\mathrm F}+\hat{\mathrm G}$ s'écrit $\hat{\mathrm A}(\mathrm F+\mathrm G)=\unsure{.}\hat{\mathrm A}\mathrm F+\hat{\mathrm A}\mathrm G$ et est donc évidente ; pour établir la relation $\handwritten{\widehat}{\mathrm F\handwritten{\cap}\mathrm G}=\hat{\mathrm F}\handwritten{\cap}\hat{\mathrm G}$ on l'écrit $\hat{\mathrm A}.(\mathrm F\handwritten{\cap}\mathrm G)=\hat{\mathrm A}.\mathrm F\handwritten{\cap}\hat{\mathrm A}.\mathrm G$, et on applique la Prop.\unsure{,} ?? du chap.~I, \handwritten{\S} ?? sur les couples plats ; même démonstration pour la relation concernant les transporteurs.

(ici encore, le rédacteur a préféré tout bloquer dans un unique

% archive sheet number: 52
\scanpage{54}
trisection de $\mathcal L$, de ``somme'' nulle.

\uline{La} \deleted{\unsure{classification}} \inserted{donnée} \uline{d'un système} \deleted{\unsure{préelliptique|préelliptiques|privilégié}} \uline{\unsure{anharmonique|anharmoniques}} \inserted{\uline{\unsure{pointés|pointé|portés}}} \uline{sur $S$} est donc équivalente \uline{fonctoriellement} \uline{à} \uline{celle d'un} \uline{\unsure{faisceaux|faisceau}} \deleted{\unsure{en}} \uline{inversible} \deleted{\illegible{1}} $\underline L$ \uline{sur} $S$, \uline{et d'une trisection} $X'$ \inserted{\uline{non ramifiée}} \uline{de} $\mathcal L=\mathrm V(\underline L)$, \uline{de} \deleted{\unsure{projectée nulle}} ``\uline{somme}'' \uline{nulle} \deleted{\unsure{dans}} \deleted{$\underline L^{\otimes3}$}. Mais on peut aussi définir les \inserted{\unsure{autres}} fonctions symétriques élémentaires \unsure{pour cette trisection}

% slanted
\margin{un \unsure{revêt\textsuperscript{mt}} \deleted{\unsure{galoisien de}} non ramifié \deleted{\unsure{groupe $G$}} d'ordre 3, et une section de}

% archive sheet number: 53
\scanpage{55}
\mspage{- 25 -}

n\textsuperscript{o} 276

\uline{couple} $(\hat{\mathrm A},\mathrm m)$ \uline{est un anneau de Zariski}.

(b) - \uline{Pour tout} $\mathrm A$\uline{-module} $\mathrm E$ \uline{de type fini}, \uline{l'application canonique}
\[
\hat{\mathrm A}\underset{\mathrm A}{\handwritten{\otimes}}\mathrm E\handwritten{\longrightarrow}\hat{\mathrm E}
\]
\uline{est bijective}.

(c) - \uline{Le couple} $(\mathrm A,\hat{\mathrm A})$ \uline{est plat}.

Soit $\mathrm u:\mathrm E\handwritten{\longrightarrow}\mathrm F$ un homomorphisme de $\mathrm A$-modules de type fini. Comme on a $\mathrm u(\mathrm m^{\mathrm p}\mathrm E)=\mathrm m^{\mathrm p}\mathrm u(\mathrm E)$, le groupe topologique $\mathrm u(\mathrm E)$ est quotient du groupe topologique $\mathrm E$ ; or la topologie $\mathrm m$-adique de $\mathrm u(\mathrm E)$ est induite par celle de $\mathrm F$ (Théorème 1) ; par conséquent, $\mathrm u$ est un \uline{homomorphisme topologique} de $\mathrm E$ dans $\mathrm F$. Il résulte de là et du lemme 1 du \handwritten{\S}1, n\textsuperscript{o}7, que le foncteur $\mathrm E\handwritten{\to}\hat{\mathrm E}$ transforme toute suite exacte de $\mathrm A$-modules \uline{de type fini} en une suite exacte.

Pour établir l'assertion (b) observons qu'on a une suite exacte
\[
\mathrm L_1\handwritten{\longrightarrow}\mathrm L_{\unsure{\mathrm o|0}}\handwritten{\longrightarrow}\mathrm E\handwritten{\longrightarrow}0
\]
avec $\mathrm L_{\mathrm o}$ et $\mathrm L_1$ libres et de type fini puisque $\mathrm A$ est noethérien. D'après ce qui précède et les propriétés élémentaires des produits tensoriels, il vient donc un diagramme commutatif de suites exactes

\diagram{\begin{tikzcd}[ampersand replacement=\&]
\hat{\mathrm A}\handwritten{\otimes}\mathrm L_1 \handwritten{\arrow[r]}\handwritten{\arrow[d]} \& \hat{\mathrm A}\handwritten{\otimes}\mathrm L_{\mathrm o} \handwritten{\arrow[r]}\handwritten{\arrow[d]} \& \hat{\mathrm A}\handwritten{\otimes}\mathrm E \handwritten{\arrow[r]}\handwritten{\arrow[d]} \& 0 \\
\hat{\mathrm L}_1 \handwritten{\arrow[r]} \& \hat{\mathrm L}_{\mathrm o} \handwritten{\arrow[r]} \& \hat{\mathrm E} \handwritten{\arrow[r]} \& 0
\end{tikzcd}}
% The nine arrows of the diagram are drawn in by hand.

De sorte que pour établir que le troisième homomorphisme vertical est bijectif il suffit d'établir que les deux premiers le sont : mais si $\mathrm E$ est isomorphe à $\mathrm A^{\mathrm n}$ il est clair que $\hat{\mathrm E}$ est isomorphe à $\hat{\mathrm A}^{\mathrm n}\unsure{.}=\hat{\mathrm A}\underset{\mathrm A}{\handwritten{\otimes}}\mathrm A^{\mathrm n}$, d'où le résultat.

Il suit de là que le foncteur $\mathrm E\handwritten{\to}\hat{\mathrm A}\underset{\mathrm A}{\handwritten{\otimes}}\mathrm E$ est exact sur la catégorie des $\mathrm A$-modules de type fini ; le $\mathrm A$-module $\hat{\mathrm A}$ est donc plat. Comme de plus l'homomorphisme $\mathrm E\handwritten{\to}\hat{\mathrm A}\underset{\mathrm A}{\handwritten{\otimes}}\mathrm E=\hat{\mathrm E}$ est injectif si $\mathrm E$ est de type fini - ceci parce que $\mathrm E$ est séparé pour la topologie \quad -adique en vertu du théorème 2-

% archive sheet number: 54
\scanpage{57}
\diagram{\begin{tikzcd}[ampersand replacement=\&, row sep=small,column sep=small]
\deleted{\Delta^2} \arrow[r,dash] \arrow[d,dash] \& \deleted{\Delta^3}\arrow[d,dash]\\
\deleted{y} \arrow[r,dash] \& \deleted{x}
\end{tikzcd}}

\[
\deleted{\unsure{x_{\xi}}}\qquad
(D_{\xi}u_{\xi})^2=4u_{\xi}^3+b_{\xi}u_{\xi}+c_{\xi}
\]
\[
\deleted{(D_{\xi}x_{\xi})^2=4\unsure{.}}\qquad
\deleted{\unsure{u_{\xi}^{2}=\Delta^3|u_\xi\in\Delta^3}}\quad
(D_{\lambda\xi}u_{\lambda\xi})^2=4u_{\lambda\xi}^3+b_{\lambda\xi}u_{\lambda\xi}+c_{\lambda\xi}
\]
\[
u_{\xi}=6\unsure{X|x}_{\xi^2}\qquad \deleted{\unsure{x}}\quad u_{\lambda\xi}=\unsure{\rho|p}u_{\xi}
\]

\diagram{\begin{tikzcd}[ampersand replacement=\&,row sep=small]
C\arrow[d]\ \big|\ \underline O\to\underline O(1) \& f_*(\underline O(n))=E_n \\
S \&
\end{tikzcd}}

\[
\left|\quad\lambda^2\rho^2(D_{\xi}u_{\xi})^2\underset{=}{=}4\rho^3u_{\xi}^3+\rho b_{\lambda\xi}u_{\xi}^{\deleted{\unsure{3}}}+c_{\lambda\unsure{\xi|\zeta}}
\right.
\]

\diagram{$\begin{array}{cccccccc}
\deleted{\unsure{0\to\underline O|\underline O=}}&
\underbrace{\phantom{x}E_1}_{\underline O_S=\Delta^0}&
\underbrace{\subset E_2}_{\deleted{\unsure{\Delta}}\,L_2}&
\underbrace{\subset E_3}_{L_3}&
\underbrace{\subset\deleted{\unsure{\subset}}E_4}_{L_2^2}&
\underbrace{\subset E_5}_{L_2L_3}&
\underbrace{\subset E_6}_{L_2^3}&\subset\cdots\quad\underbrace{\subset E_n}\\
&&\substack{\text{\unsure{somme}}\\\text{des valeurs}\\\text{en }\unsure{\omega_1,\omega_2,\omega_3}\\\text{est nulle}}&
\substack{\Vert\\\Delta L_2}&
\deleted{\substack{\vert\!\vert\!\vert\\\Delta L_3\\\unsure{L_2^2}}}&
\deleted{\substack{\vert\!\vert\!\vert\\\Delta^3 L_2}}&
\deleted{\substack{\vert\!\vert\!\vert\\L_3^2\\\unsure{\mathcal L}^{\prime\prime\prime}}}&\\
&\multicolumn{2}{c}{\underbrace{\hspace{4em}}_{\text{\unsure{paires|pairs}}}}&&&&&
\end{array}$}

\[
\deleted{\unsure{\xi}}\quad y^{2}=4X^3+bx+c\qquad
\frac{y}{\deleted{x}}=D
\]

\[
(D_{\xi}u_{\xi})^2=4\frac{\rho}{\lambda^2}u_{\xi}^3+
\frac{\unsure{1|b}}{\rho\lambda^2}b_{\lambda\xi}u_{\unsure{\xi|\zeta}}+
\frac{c_{\lambda\xi}}{\lambda^2\rho^2}\ .
\]

\[
\rho=\lambda^2\qquad
\left\{\begin{aligned}
u_{\lambda\xi}&=\lambda^2u_\xi\\
b_{\lambda\xi}&=\lambda^4 b_\xi\\
c_{\lambda\unsure{\xi|\zeta}}&=\lambda^6c_\xi
\end{aligned}\right.
\qquad
\begin{aligned}
u_\xi&\in L_2\simeq\Delta^2\\
b_\xi&\in\deleted{\unsure{O}}\Delta^0\\
c_\xi&\in\Delta^0
\end{aligned}
\]

\[
\Delta^4\longrightarrow\underline O_S\qquad
\Delta^6\unsure{\mathbin{-}|\longrightarrow}\underline O_S\qquad
\Delta^2\longrightarrow\Delta^2\qquad
\begin{aligned}
\Delta^{\unsure{n|12}}&\simeq\mathcal L^{18}\\
\Delta^2&\simeq\mathcal L^3
\end{aligned}
\]

\[
\left\{\begin{aligned}
b&\in\Delta^{-4}\\
c&\in\Delta^{-6}\\
u&\quad\text{scalaire}
\end{aligned}\right.
\qquad
u_\xi=\xi^2
\]

\figure{\begin{tikzpicture}[x=1cm,y=1cm]
\draw (0,0) -- (1.25,1.15) -- (4.0,1.08) -- (7.25,1.20);
\node at (2.28,.59) {$\xi^2$};
\end{tikzpicture}}
% archive sheet number: 56
\scanpage{58}
\[
\deleted{\unsure{X_1}} \qquad\qquad B\qquad A
\]

\[
X\qquad X\times_S X\supset R
\]

\diagram{\begin{tikzcd}[ampersand replacement=\&,column sep=small]
\& X \arrow[r] \& Y\\
X\times_{\unsure{Y|y}}X\arrow[ur]\arrow[r] \& X\arrow[ur] \&
\end{tikzcd}}

\[
\begin{aligned}
&&x_i\\
S&& B\underset{A}{\otimes}B
\end{aligned}
\]
\[
\unsure{1|\mid}
\]

\diagram{\begin{tikzcd}[ampersand replacement=\&, row sep=small]
B\arrow[d,dash] \& A\arrow[r,dash] \& B\otimes_A B\\
A \& \& B
\end{tikzcd}}

\[
x_i\otimes 1\quad\text{base de}\quad B\otimes B\quad\text{sur}\quad 1\otimes B
\]
\margin{$/$ \unsure{divisant|diviser} par $\unsure{m|\mathfrak m}\subset A$}

\[
\dot x_i\otimes 1\qquad\mathord{\text{---}}\quad L\otimes_{\deleted{\unsure{A}}k}L\quad\text{sur}\quad1\otimes L
\]

$x_i$ base de $L$ sur $k$

$x_i$ \unsure{générateurs} de $B$ sur $A$ par Nakayama.

\[
\sum\lambda_i x_i=0\qquad\lambda_i\in A
\]

\[
\deleted{\sum\lambda_i x_i}\left(\sum\lambda_i x_i\right)\otimes_A1
=\sum(x_i\otimes1)(1\otimes\lambda_i)=0
\]

\[
1\otimes\lambda_i=0\quad\text{pour }\text{tt}\quad i
\]

\unsure{Utilisant|Appliquant} $B\otimes_A B\unsure{\mathbin{-}|\to}B$, on trouve.

\[
\deleted{\lambda_i\otimes1=0\qquad1\otimes\lambda_i=\lambda_i\otimes1}
\]

\[
\lambda_i=0\quad\text{dans}\quad B
\]
% archive sheet number: 57
\scanpage{59}
Théorie transcendante

(formules classiques)

% archive sheet number: 58
\scanpage{60}
\uline{fonction elliptique}. C'est une fonction méromorphe dans le plan à distance finie, et admettant deux périodes \unsure{distinctes|distincte} $\omega_1$ et $\omega_2$. En particulier, une constante est une fonction elliptique. Si $f(z)$ n'est pas une constante, le rapport de $\omega_1$ et $\omega_2$ n'est \deleted{imaginaire} \inserted{pas réel}.

\inserted{\unsure{Définition} de la période ($\omega_1$ et $\omega_2$ \unsure{donnés|donnent} une fois pour \unsure{ttes}) La donnée d'une fonction elliptique n'entraîne \uline{pas} la donnée de $\omega_1$ et $\omega_2$. Mais il faut remarquer que}

Pavage correspondant du plan. On peut définir un réseau de parallélogrammes \unsure{correspondants|équipollents} aux périodes $\omega_1$ et $\omega_2$ et à un sommet $a$ arbitraire. Ils sont \unsure{déterminé|déterminés} par leurs sommets de la forme $a+n_1\omega_1+n_{\unsure{1|2}}\omega_2$ ; tout point qui n'est pas de la forme $a+n\omega_1+\lambda\omega_1$, ou $a+\lambda\omega_1+n\omega_2$ ($\unsure{b|\lambda}$ et $n$ réels, $n$ entier) appartient : 1 et 1 seul de ces parallélogrammes, \inserted{\unsure{il lui est intérieur|il lui est int.}} \unsure{le pt|tout point} qui est de la forme précédente, mais non de la forme $a+n_1\omega_1\unsure{-|+}n_1\omega_2=a+\omega$, \deleted{est} appartient à deux et deux seult de ces par. \deleted{\unsure{en}} \inserted{il appartient aux côtés, non au sommet} enfin, un point de la forme $a\unsure{-|+}n_1\omega_1\unsure{-|+}n_1\omega_2$ appartient à 4 et 4 seult de tels par, il appartient à leurs sommets.

Si on se donne un parallélogramme de périodes, tout point est l'``homologue'' soit d'un pt intérieur, soit de deux p\textsuperscript{ts} sur les côtés, et non les sommets, soit des 4 sommets. \inserted{En particulier, deux p\textsuperscript{ts} \uline{int} à un $\overline{\mathrm m}$ par. de périodes ne sont pas homologues} \deleted{Si on se donne la fonction}

On peut trouver des par. des périodes, tels que sur leur côté ne se trouve aucun zéro, pôle etc. ; Dans ce cas, la \deleted{\unsure{nb}} \inserted{liste} des zéros (pôles) avec leur ordre de multiplicité, est \unsure{indépendant|indépendante} du par. considéré. Il en est ainsi en tte généralité, si on \unsure{cpte|compte} \deleted{une seule fois} \deleted{\unsure{disting}}é les zéros de tout groupe de zéros homologues (\deleted{\unsure{situé|situés}} ils est situés sur les sommets - auquel cas ils sont 4, on les divise par 4 -- autrement on les divise par deux)

% archive sheet number: 59
\scanpage{61}
L'ens. des val. prises par $f(z)$ dans un par. des périodes est égal à l'ensemble des valeurs prises dans tout le plan ;

\margin{$\times$}Si une fonction elliptique n'a pas de pole, elle est constante.

Remarquer que la somme, le produit, le quotient, la dérivée, de fonction elliptique de $\unsure{\overline{\mathrm m}^{\mathrm s}|=^{\mathrm s}}$ périodes ont $\overline{\mathrm m}$ périodes.

\margin{$\uparrow$}\uline{Théorème $\overline{\mathrm{II}}$} La somme des résidus d'une fonction elliptique dans un par. de périodes est nulle.

Donc, une fonction elliptique admet dans un par. de périodes au moins deux poles (chacun compté avec son ordre de multiplicité)

\uline{Théorème $\overline{\mathrm{III}}$} Le nombre des zéros et le nb des poles d'une fonction elliptique dans un parallélogramme de période est le même (en supposant tjrs que sur les bords --) ; il en est de même du nb de fois qu'il pd \deleted{la valeur} \inserted{une valeur} donnée (sous les $\overline{\mathrm m}$ réserves)

Ce nombre est appellé ordre de la fonction elliptique.

\margin{$\downarrow$}\uline{Théorème $\overline{\deleted{\unsure{\mathrm{III}|\mathrm{II}|\mathrm I}}}$} Deux fonctions elliptiques de $\overline{\mathrm m}$ périodes, qui ont mêmes zéros et même poles, avec le $\overline{\mathrm m}$ ordre de multiplicité, sont égales à un facteur constant près ; Elles sont égales à une constante additive près, si elles ont $\overline{\mathrm m}$ périodes, et $\overline{\mathrm m}^{\mathrm s}$ poles, avec $\overline{\mathrm m}^{\mathrm s}$ parties principales.

\uline{Théorème $\overline{\mathrm{IV}}$} la somme des zéros d'une fonction \deleted{\unsure{holo|toute}} elliptique ds un parallélogramme de périodes égale la somme de ses poles, à une période près.

\uline{Théorème $\overline{\mathrm V}$}. \unsure{Ent|Entre} deux fonctions elliptiques de $\overline{\mathrm m}$ période Existe une relation algébrique

% archive sheet number: 60
\scanpage{62}
\uline{La fonction $\unsure{p|\wp}z$}

Définition intrinsèque : c'est une fonction elliptique d'ordre 2, admettant l'origine pour pole double, \deleted{\unsure{y ayant pour pa}} \inserted{la fonction} $pz-\dfrac{1}{z^2}$ s'annulant à l'origine. Si $\omega_1$ et $\omega_2$ sont donnés, cela détermine sans plus cette fonction elliptique.

Elle est paire [car $\varphi(z)-\varphi(-z)$ est elliptique, n'a pas de pole, s'annule à l'origine] \deleted{et vérifie l'équation} \inserted{et est donc de la forme}
\[
\deleted{\unsure{p^{\prime2}z}}\qquad
pz=\frac1{z^2}+c_1z^{\unsure{2|3}}+c_{\unsure{1|2}}z^4\mathord{\text{---}}
\]

Elle vérifie l'équation
\[
p^{\prime 2}z=4p^3z-g_2pz-g_3\qquad
g_{\unsure{1|2}}=20c_1\qquad g_3=28c_2
\]

C'est donc la fonction inverse de
\[
z=\int_\infty^Z\frac{dZ}{\sqrt{4Z^3-g_2Z-g_3}}
\]

On en déduit que si on pose
\[
p\omega_1=e_1\qquad p\omega_2=e_2\qquad p\omega_3=e_3\ ,
\]
$e_1,e_2,e_3$ sont les racines de l'équation $4Z^3-g_2Z-g_3=0$ ; on a donc $p'\omega_1=p'\omega_2=p'\omega_3=0$

On a \deleted{$p'\omega_1=\unsure{e_1}$}
\[
\begin{aligned}
p''z&=6p^{\unsure{2|\prime}}z-\frac{g_2}{2}\\
p^{\unsure{\prime\prime\prime|\prime\prime}}z&=12pz\,p'z\\
&\mathord{-}\quad\mathord{-}\quad\mathord{-}\quad\mathord{-}
\end{aligned}
\]

$p^{\unsure{(2n)|(2n-1)}}z$ s'exprime en fonction de $pz$ par un polynome de degré $n+1$ ; $p^{2n+1)}z$ par le produit par $p'z$ d'un polynome en $pz$ de degré $n$.

\inserted{Développ\textsuperscript{mt} en série double}

\uline{La fonction $\zeta z$}. S'introduit par l'intégration des fonctions elliptiques.

La fonction $\int_0^z(pz-\frac1{z^2})dz$ est uniforme, en lui ajoutant $\frac1z$ on obtient une primitive de $\unsure{pz|\wp z}$. On pose
\[
\zeta z=\frac1z-\int_0^z(pz-\frac1{z^2})dz\qquad\text{On a }\quad\zeta'z=-pz
\]
% archive sheet number: 61
\scanpage{63}
$\zeta z$ est une fonction impaire, \deleted{\unsure{non périodique}. On a} \inserted{admettant les \unsure{nbs} $\omega$ pour pole simple (avec 1 comme résidu)}
\[
\zeta(z+2\omega_1)-\zeta z=c^{\mathrm{te}}=2\zeta\omega_1=2\eta_1\qquad
\text{on définit de }\unsure{\overline{\mathrm m}|=}\ \eta_2\text{ et }\eta_3
\]

On établit par une intégration de $\int z\,\unsure{p|\wp}z\,dz$ sur un contour convenable
\[
\omega_2\eta_1-\omega_1\eta_2=\frac{\pi i}{2}\qquad\text{etc. ; d'ailleurs }\quad\eta_1+\eta_2+\eta_3=0
\]

La fonction $\zeta z$ permet de calculer l'intégrale $\displaystyle\int\frac{Z\,dZ}{\sqrt{4Z^3-g_2Z-g_3}}$ on pose $Z=pz$ d'où $\displaystyle\int=\int pz\,dz=-\zeta z+c^{\mathrm{te}}$

\deleted{D'une manière \unsure{précise|générale},}

\uline{Formule de décomposition de Hermite}. Soit $f(z)$ une fonction elliptique de périodes $2\omega_1$ et $2\omega_2$ ; $a,b,\mathord{--}$ ses poles, distincts à une période près deux à deux,
\[
\frac{A_\alpha}{(z-a)^\alpha}+\mathord{--}+\frac{A_1}{z-a}\ ;\qquad
\frac{B_\beta}{(z-b)^\beta}\mathord{--}+\frac{\unsure{b|B}_{\unsure{1|2}}}{z-b}\ ;\ \mathord{---}
\]
les parties principales de ses poles ; On a alors
\[
\begin{aligned}
f(z)=c^{\mathrm{te}}+{}&\left\{A_1\zeta(z-a)+A_2p(z-a)-\frac{A_3}{2!}p'(z-a)+\mathord{--}+(-1)^\alpha\frac{A_\alpha}{(\alpha-1)}p^{(\alpha-2)}(z-a)\right\}\\
+{}&\{B_1\zeta(z-b)+\mathord{---}\qquad\qquad\}\\
&\mathord{-}\quad\mathord{-}\quad\mathord{-}\quad\mathord{-}\quad\mathord{-}\quad\mathord{-}\quad\mathord{-}
\end{aligned}
\]

\fbox{\uline{fonction $\sigma z$}} Nous savons intégrer toutes les fonctions elliptiques au moyen de fonctions connues, si nous savons intégrer $\zeta z$. Or, son intégrale est complexe, \deleted{et} admet \unsure{tous} les points $\omega$ pour points logarithmiques. Cela suggère de chercher $\sigma z$ tel que $\int\zeta z\,dz=\log\sigma z$\qquad $\sigma z$ s'annulant aux points $\omega$\unsure{--}

% archive sheet number: 62
\scanpage{64}
On peut montrer directement que $\sigma z=e^{\int\zeta z\,dz}$ est uniforme, \deleted{est donc} \inserted{est même} une fonction entière ; \deleted{il est} on peut aussi procéder directement.

On a \deleted{en effet} $\displaystyle\unsure{p|\wp}z=\frac1{z^2}+\sum{\sum}'\left[\frac1{(z-\omega)^2}-\frac1{\omega^2}\right]$ la \unsure{somme|série} étant unift convergente pour $z$ borné, à distance $>\varepsilon$ des $\omega$. Les dérivées successives \unsure{s'obtient|s'obtiennent} par dérivation terme à terme, la convergence est comme précédemment.

Notons que l'on obtient ainsi les coefficients du développement de $pz$
\[
c_1=3\sum{\sum}'\frac1{\omega^{\unsure{2|4}}}\qquad c_2=5\sum{\sum}'\frac1{\omega^6}\qquad\mathord{---}
\]

On a
\[
\zeta z=\frac1z+\sum{\sum}'\left[\frac1{z-\omega}+\frac1\omega+\frac{z}{\deleted{\unsure{2}}\omega^2}\right]
\]

Posons enfin
\[
\sigma z=z\prod{\prod}'\left(1-\frac z\omega\right)e^{\frac z\omega+\frac{z^2}{2\omega^2}}
\]

Le produit converge \deleted{et} uniformément \inserted{et absol\textsuperscript{t}} si $z$ est borné ; $\sigma z$ est une fonction entière, qui s'annule pour les $z=\omega$ et pas d'autres points. Elle a été construite de sorte à avoir
\[
\frac{\sigma'z}{\sigma z}=\zeta z
\]

(À remarquer que, si on voulait construire à priori une fonction entière s'annulant aux points $\omega$, aussi simple que possible, c'est $\sigma z$ qu'on aurait trouvé)

$\sigma z$ est une fonction paire. On a
\[
\sigma(z+2\omega_1)=-e^{2\eta_1(z+\omega_1)}\sigma z\qquad
\sigma(z+\omega_2)=-e^{2\eta_2(z+\omega_{\unsure{2|1}})}\sigma z
\]
de $\unsure{\overline{\mathrm m}|=}$ pour $\omega_3$

% archive sheet number: 63
\scanpage{65}
\uline{formule de Jacobi} Soient $a,b,\mathord{\text{---}}\ l$ les poles \deleted{d'une} \inserted{d'une} fonction elliptique de \unsure{pér|pés} $2\omega_1$ et $2\omega_2$ (distincts à \unsure{$2\omega$|$1\omega$} près) $a_1\ b_1\ \mathord{\text{---}}\ l_1$ les zéros ; on a
\[
a+b\mathord{---}+l=a_1\unsure{+|-}b_1\mathord{---}+l_1+\Omega\ .\quad \Omega\ \text{est une}
\]

On a alors
\[
f(z)=c^{\mathrm{te}}\times\frac{\sigma(z-a_1)\sigma(\unsure{z|B}-b_1)\mathord{--}\sigma(z-l_1-\Omega)}
{\sigma(z-a)\sigma(z-b)\mathord{---}\sigma(z-l)}
\]
soit
\[
f(z)=c^{\mathrm{te}}\times\frac{\sigma(z-a_1)\sigma(z-b_1)\mathord{-}\sigma(z-l_1)}
{\sigma(z-a)\sigma(z-b)\mathord{---}\sigma(z-l+\Omega)}
\]
soit \unsure{aussi|enfin|encore}, enfin
\[
f(z)=c^{\mathrm{te}}\times e^{2(n_1\eta_1+n_2\eta_2)z}
\frac{\sigma(z-a_1)\mathord{--}\sigma(z-l_1)}{\sigma(z-a)\mathord{--}\sigma(z-l)}
\]
en posant $\Omega=2n_1\omega_1+2n_2\omega_2$

\uline{formule d'addition}

\[
\begin{aligned}
\unsure{p|\wp}u-p\unsure{v|w}&=-\frac{\sigma(u+v)\sigma(u-v)}{\sigma^2u\,\sigma^2v}\\
\zeta(u+v)&=\zeta u+\zeta v+\frac12\,\frac{p'u-p'v}{pu-pv}\\
p(u+\unsure{u|v})+pu+pv&=\frac14\left(\frac{p'u-p'v}{pu-pv}\right)^2=[\zeta(u+v)-\zeta u-\zeta v]^2
\end{aligned}
\]
% archive sheet number: 64
\scanpage{66}
\begingroup\makeatletter\tagsleft@true\makeatother
\uline{Fonction de Legendre $\mathrm{sn}\,z$}

La forme $\displaystyle\int\frac{dZ}{\sqrt{(1-Z^2)(1-k^2 Z^2)}}$ est tentante, car elle paraît très symétrique, et rappelle des intégrales élémentaires \unsure{liées|liés} aux fonctions \unsure{trigonométrique|trigonométriques} et \unsure{hyperbolique|hyperboliques} ($k=0,\ k=1$).

Appelons $\mathrm{sn}\,z$ l'intégrale de l'équation
\[
\left(\frac{dZ}{dz}\right)^2=(1-Z^2)(1-k^2Z^2)
\]
qui pd la valeur zéro au point $z=0$, avec $Z'(0)=+1$. L'unicité de cette solution se montre de diverses façons.

1\textsuperscript{o} En faisant un changement de variable convenable, on obtient
\[
Z=\frac{12\unsure{p|\wp}(z-z_0)-(5-k^2)}{12p(z-z_0)-(5k^2-1)}
\qquad
\begin{aligned}
e_1&=\tfrac16(1\unsure{+|\cdot}k^2)\\
e_2&=-\tfrac1{12}(k^2+6k+1)\\
e_3&=-\tfrac1{12}(k^2-6k+1)
\end{aligned}
\]

L'intégrale particulière est
\[
\mathrm{sn}\,z=\frac{12p(z-z_0)-(5-k^2)}{12p(z-z_0)-(5k^2-1)}
\qquad\text{avec}\quad pz_0=\frac{5-k^2}{12}\quad p'z_0=-\frac{1-k^2}{2}\ .\tag*{(1)}
\]

Ces formules sont de forme trop lourde, eu égard à d'autres plus simples. Mais on voit l'existence de $\mathrm{sn}\,z$, qui est une fonction elliptique dont les périodes s'expriment par des intégrales.

\endgroup
% archive sheet number: 65
\scanpage{67}
\begingroup\makeatletter\tagsleft@true\makeatother
\uline{2°} On obtient $Z^2=\dfrac{1}{k^2}\bigl[\unsure{p|\wp}(z-z_0)-e_1\bigr]$
\[
\left\{\begin{aligned}
e_1&=-\frac{1+k^2}{\deleted{\unsure{3}}3}\\
e_3&=\frac{2k^2-1}{3}\\
e_2&=\frac{2-k^2}{3}
\end{aligned}\right. \tag*{(2)}
\]

La solution envisagée est
\[
\boxed{\operatorname{sn}^2 z=\frac{1}{k^2}\bigl[p(z-\omega_1)-e_1\bigr]} \tag*{(3)}
\]

\uline{3°} Avec la même fonction elliptique, on a une solution générale $Z^2=\dfrac{1}{p(z-z_0)-e_1}$ et la solution particulière est
\[
\boxed{\operatorname{sn}^2 z=\frac{1}{pz-e_1}} \tag*{(4)}
\]

On montrerait facilement l'uniformité de \deleted{$Z(z)$} $\operatorname{sn}z$ (les racines de $pz-e_1=0$ étant doubles)

$\ast\ast\ast$ Ces deux formules impliquent une identité :
\[
(pz-e_1)(p(z-\omega_1)-e_1)=k^2
\]
facile à vérifier, en passant à la fonction $\sigma$, ou en montrant que le premier mb qui est elliptique, n'a pas de pole, et est donc constant -- ceci est vrai pour \deleted{tte} fonction $pz$ -- On trouve d'ailleurs incidemment la formule $\dfrac{e^{2\eta_1\omega_1}}{\sigma^{\unsure{4|2}}\omega_1}=3e_1^2-\dfrac{g_2}{4}$, \deleted{qui permet le calcul de $\eta_1$ et \unsure{à|de} fort. $\eta_2$ et $\eta_3$}

Dans le cas actuel, on a $k^2=\dfrac{e^{2\eta_1\omega_1}}{\sigma^4\omega_1}=3e_1^2-\dfrac{g_2}{4}$ $\,-\ast\ast\ast$

Toujours est il que nous avons $\operatorname{sn}^2(z+\omega_1)\operatorname{sn}^2z=\dfrac{\deleted{\unsure{1}}1}{k^2}$ ; si nous choisissons $\omega_1$ et $k$ convenablement nous aurons
\[
\boxed{\begin{aligned}
&\operatorname{sn}(z+\omega_1)\operatorname{sn}z=\frac{1}{k}\\[-.5ex]
&\scriptstyle\unsure{=}
\end{aligned}} \tag*{(5)}
\]
\endgroup
% archive sheet number: 66
